\makeatletter
\@ifundefined{keepfig}{\newenvironment{keepfig}{\par\edef\keepfig@pi{\the\parindent}\noindent
  \begin{minipage}{\linewidth}\parindent=\keepfig@pi\relax}{\end{minipage}\par}}{}
\makeatother
\section*{Упражнения}
\addcontentsline{toc}{section}{Упражнения}

% ---------------------------------------------------------------- 3-1
\begin{exercise}[3-1]
Пусть $F$~--- поле характеристики $0$. Покажите, что $F(X^2)\cap F(X^2-X)=F$ (пересечение берётся внутри $F(X)$).
[Указание: найдите автоморфизмы $\sigma$ и $\tau$ поля $F(X)$ порядка $2$, оставляющие на месте соответственно $F(X^2)$ и
$F(X^2-X)$, и покажите, что $\sigma\tau$ имеет бесконечный порядок.]
\end{exercise}

\noindent\emph{Идея.} Поле $E=F(X^2)\cap F(X^2-X)$ неподвижно относительно обоих автоморфизмов, а значит, и относительно
порождённой ими \emph{бесконечной} группы. С другой стороны, над подполем, содержащим хотя бы одну непостоянную
рациональную функцию, поле $F(X)$ конечно, а у конечного расширения лишь конечное число автоморфизмов (\rf{2.14}).
Значит, в $E$ нет непостоянных функций. Нужна лемма \rf{3.1a} (она верна в любой характеристике): \emph{если
$u=a/b\in F(X)\smallsetminus F$, где $a,b\in F[X]$ взаимно просты, то $X$ алгебраичен над $F(u)$ и
$[F(X):F(u)]=\deg u:=\max(\deg a,\deg b)$.} Например, для $u=\frac{X^2+X+1}{X+1}=X+\frac1{X+1}$ элемент $X$~--- корень
многочлена $T^2+T+1-u(T+1)=T^2+(1-u)T+(1-u)$, и $[F(X):F(u)]=2$. В частности, $F(u)=F(X)$ ровно тогда, когда $\deg u=1$
(\rf{3.1b}(a), \rf{9.25}); а по теореме Люрота (\rf{9.19}) каждое промежуточное поле $F\subsetneq K\subset F(X)$ имеет вид
$F(u)$~--- но для упражнения это не нужно.

\medskip\noindent\emph{Решение.} Подстановка $X\mapsto cX+d$ ($c,d\in F$, $c\ne0$) задаёт $F$-автоморфизм кольца $F[X]$ (обратный~---
подстановка $X\mapsto(X-d)/c$), а значит, и его поля частных $F(X)$: $f(X)\mapsto f(cX+d)$. Пусть
\[
\sigma\colon X\mapsto -X,\qquad \tau\colon X\mapsto 1-X .
\]
Тогда $\sigma^2=\tau^2=1$, $\sigma,\tau\ne1$ (так как $2\ne0$ в $F$), и
\[
\sigma(X^2)=X^2,\qquad \tau(X^2-X)=(1-X)^2-(1-X)=X^2-X .
\]
Автоморфизм, оставляющий на месте образующую поля, оставляет на месте и всё поле; значит, $\sigma$ тождествен на $F(X^2)$, а
$\tau$~--- на $F(X^2-X)$. Более того, $F(X^2)=F(X)^{\langle\sigma\rangle}$: по \rf{3.1a} $[F(X):F(X^2)]=2$, по теореме Артина
(\rf{3.11}) $[F(X):F(X)^{\langle\sigma\rangle}]=2$, и $F(X^2)\subset F(X)^{\langle\sigma\rangle}$. Так же
$F(X^2-X)=F(X)^{\langle\tau\rangle}$. Поэтому
\[
E:=F(X^2)\cap F(X^2-X)=F(X)^{\langle\sigma,\tau\rangle}
\]
(элемент, неподвижный относительно $\sigma$ и $\tau$, неподвижен относительно всех их произведений; ср.\ \rf{3.18}(a) и
\rf{3.18a}).

Положим $\rho=\sigma\tau$ (сначала $\tau$, потом $\sigma$): $\rho(X)=\sigma(1-X)=1+X$. По индукции $\rho^m(X)=X+m$, и при
$m\ge1$ имеем $\rho^m\ne1$, так как $m\cdot1\ne0$ в поле характеристики $0$. Итак, группа $\langle\sigma,\tau\rangle\subset\Aut(F(X)/E)$
бесконечна.

Предположим, что $E\ne F$, и возьмём $u\in E\smallsetminus F$. Тогда $F(u)\subset E\subset F(X)$, и по \rf{1.20} и
\rf{3.1a}
\[
[F(X):E]\le[F(X):F(u)]=\deg u<\infty .
\]
Но по \rf{2.14} число $E$-автоморфизмов конечного расширения $F(X)/E$ не больше $[F(X):E]$~--- противоречие с
бесконечностью $\Aut(F(X)/E)$. Значит, $E=F$. \qed

\begin{center}
\begin{keepfig}
\centering
\begin{tikzpicture}[every node/.style={font=\small}]
\node (FX) at (0,2.6) {$F(X)$};
\node[cold] (A) at (-2.2,1.3) {$F(X^2)=F(X)^{\langle\sigma\rangle}$};
\node[cold] (B) at (2.2,1.3) {$F(X^2-X)=F(X)^{\langle\tau\rangle}$};
\node[cnew] (E) at (0,0) {$E=F(X)^{\langle\sigma,\tau\rangle}=F$};

\draw (FX)--(A) node[midway,above left=-1pt] {$2$};
\draw (FX)--(B) node[midway,above right=-1pt] {$2$};
\draw (A)--(E); \draw (B)--(E);
\node[cnew,right,align=left] at (2.0,0) {$\langle\sigma,\tau\rangle\ni\rho=\sigma\tau\colon X\mapsto X+1$\\ (бесконечный порядок)};
\end{tikzpicture}
\figcap{Каждое из подполей $F(X^2)$, $F(X^2-X)$ имеет степень $2$, но их пересечение~--- неподвижное поле бесконечной
группы $\langle\sigma,\tau\rangle$~--- имеет бесконечный индекс и равно $F$.}
\end{keepfig}
\end{center}

\begin{remark}
(1) В рассуждении важно, что лемма \rf{3.1a} относится к \emph{любой} непостоянной рациональной функции, а не только к
многочленам: промежуточное поле $F\subsetneq K\subset F(X)$ может вовсе не содержать непостоянных многочленов. Например,
$K=F(X+X^{-1})$ неподвижно относительно автоморфизма $X\mapsto X^{-1}$, а многочлен $h$ степени $n\ge1$ не может удовлетворять
$h(X^{-1})=h(X)$: умножив на $X^n$, получили бы слева многочлен степени $\le n$, а справа~--- степени $2n$.

(2) \emph{Прямое доказательство} (без степеней расширений). Каждый $u\in E$ неподвижен относительно $\rho$, т.\,е.\
$u(X+1)=u(X)$. Запишем $u=a/b$, $\gcd(a,b)=1$, $b$ унитарен. Тогда $a(X+1)b(X)=a(X)b(X+1)$; многочлены $a(X+1)$ и
$b(X+1)$ взаимно просты (подстановка~--- автоморфизм $F[X]$), поэтому $b(X+1)\mid b(X)$, а так как оба многочлена
унитарны одной степени, $b(X+1)=b(X)$. Отсюда $b(n)=b(0)$ для всех $n\in\Z$, и многочлен $b-b(0)$ имеет бесконечно много
корней $n\cdot1$ (они различны в характеристике $0$), т.\,е.\ $b=1$. Тогда $a(X+1)=a(X)$, и так же $a\in F$. Итак, уже
$F(X)^{\langle\rho\rangle}=F$.

(3) \emph{В характеристике $p>0$ утверждение неверно}: там $\rho^p=1$, и группа $\langle\sigma,\tau\rangle$ конечна. Положим
$u=(X^p-X)^2$, $Y=X^2-X$. При $p=2$ имеем $u=X^4+X^2=Y^2$. При нечётном $p$ многочлен $X^p-X=X(X^{p-1}-1)$ нечётен, так
что $u=X^2(X^{p-1}-1)^2\in F[X^2]$; с другой стороны, $X^p-X=\prod_{c\in\F_p}(X-c)$, и при разбиении $\F_p\smallsetminus\{\frac12\}$
на пары $\{c,1-c\}$ получаем $(X-c)(X-1+c)=Y+c-c^2$ и $(X-\frac12)^2=Y+\frac14$, откуда
\[
u=\Bigl(Y+\tfrac14\Bigr)\prod_{\{c,1-c\}}\bigl(Y+c-c^2\bigr)^2\in F[X^2-X].
\]
Итак, $u\in E$ и $u\notin F$. При нечётном $p$ даже $E=F(u)$: группа $\langle\sigma,\tau\rangle$ состоит из $2p$ подстановок
$X\mapsto\pm X+k$, $k\in\F_p$ (диэдральная группа порядка $2p$), поэтому по \rf{3.11} $[F(X):E]=2p=\deg u=[F(X):F(u)]$, а
$F(u)\subset E$.
\end{remark}

% ---------------------------------------------------------------- 3-2
\begin{exercise}[3-2]
Пусть $p$~--- нечётное простое число, $\zeta$~--- первообразный корень $p$-й степени из $1$ в $\CC$, $E=\Q[\zeta]$ и
$G=\Gal(E/\Q)$; таким образом, $G=(\Z/p\Z)^\times$. Пусть $H$~--- подгруппа индекса $2$ в $G$. Положим
$\alpha=\sum_{i\in H}\zeta^i$ и $\beta=\sum_{i\in G\smallsetminus H}\zeta^i$. Покажите, что
\begin{enumerate}[label=(\alph*)]
\item $\alpha$ и $\beta$ неподвижны относительно $H$;
\item если $\sigma\in G\smallsetminus H$, то $\sigma\alpha=\beta$ и $\sigma\beta=\alpha$.
\end{enumerate}
Таким образом, $\alpha$ и $\beta$~--- корни многочлена $X^2+X+\alpha\beta\in\Q[X]$. Вычислите $\alpha\beta$ (или
$\alpha-\beta$) и покажите, что неподвижное поле $H$ равно $\Q[\sqrt p]$ при $p\equiv1\pmod4$ и $\Q[\sqrt{-p}]$ при
$p\equiv3\pmod4$.
\end{exercise}

\noindent\emph{Решение.} \emph{Группа $G$ и подгруппа $H$.} Многочлен $\Phi_p=X^{p-1}+\dots+X+1$ неприводим (\rf{1.16a}(c)),
поэтому это минимальный многочлен $\zeta$, и $[E:\Q]=p-1$. Его корни $\zeta,\zeta^2,\dots,\zeta^{p-1}$ лежат в $E$, так что
$E$~--- поле разложения $\Phi_p$, расширение Галуа (\rf{3.10}), и $|G|=p-1$. Каждый $\sigma\in G$ переводит $\zeta$ в корень
$\Phi_p$, т.\,е.\ $\sigma(\zeta)=\zeta^i$ с $i\in(\Z/p\Z)^\times$; обозначим такой автоморфизм $\sigma_i$, так что
$\sigma_i(\zeta^k)=\zeta^{ik}$. Отображение $\sigma_i\mapsto i$~--- инъективный гомоморфизм $G\to(\Z/p\Z)^\times$
($\sigma_i\sigma_j(\zeta)=\sigma_i(\zeta^j)=\zeta^{ij}$), а порядки групп равны; это изоморфизм (ср.\ \rf{5.8}), и мы
отождествляем $G=(\Z/p\Z)^\times=\F_p^\times$.

Группа $\F_p^\times$ циклическая (упражнение 1-3), поэтому подгруппа индекса $2$ в ней единственна. Это подгруппа
\emph{квадратов}: возведение в квадрат~--- гомоморфизм с ядром $\{\pm1\}$, так что квадратов ровно $(p-1)/2$. Итак, $H$~---
квадратичные вычеты по модулю $p$, $G\smallsetminus H$~--- невычеты; выбирать образующую $G$ (явный изоморфизм
$\F_p^\times\cong\Z/(p-1)\Z$) для этого не нужно. Например, $H=\{1,2,4\}$ при $p=7$, $H=\{1,3,4,9,10,12\}$ при $p=13$.
Нам понадобится ещё один факт: \emph{$-1\in H\iff p\equiv1\pmod4$}. Действительно, $-1$~--- единственный элемент порядка
$2$ в циклической группе $\F_p^\times$; если $-1=y^2$, то $y$ имеет порядок $4$, и обратно, квадрат элемента порядка $4$
имеет порядок $2$; а элемент порядка $4$ в циклической группе порядка $p-1$ существует тогда и только тогда, когда $4\mid p-1$.

\emph{(a) и (b).} $\sigma_i(\alpha)=\sum_{h\in H}\zeta^{ih}=\sum_{h'\in iH}\zeta^{h'}$, и так же для $\beta$. Если $i\in H$, то
$iH=H$ и $i(G\smallsetminus H)=G\smallsetminus H$, поэтому $\sigma_i\alpha=\alpha$, $\sigma_i\beta=\beta$. Если $i\notin H$, то
$iH=G\smallsetminus H$ и $i(G\smallsetminus H)=H$ (два смежных класса подгруппы индекса $2$), поэтому $\sigma_i\alpha=\beta$,
$\sigma_i\beta=\alpha$. (При $i\in H$ отображение $h\mapsto ih$~--- перестановка множества $H$, но, вообще говоря, не
цикл: её циклы~--- смежные классы $h\langle i\rangle$ длины $\operatorname{ord}i$; одним циклом она будет, лишь если $i$
порождает $H$.)

Следовательно, $\alpha+\beta$ и $\alpha\beta$ неподвижны относительно всей группы $G$ и лежат в $E^G=\Q$. Далее,
$\alpha+\beta=\zeta+\zeta^2+\dots+\zeta^{p-1}=\Phi_p(\zeta)-1=-1$, и по теореме Виета $\alpha$, $\beta$~--- корни
$X^2+X+\alpha\beta$.

\emph{Вычисление $\alpha\beta$.} Раскрыв скобки, получаем
\[
\alpha\beta=\sum_{i\in H,\ j\notin H}\zeta^{i+j}=\sum_{k\in\Z/p\Z}N(k)\,\zeta^k,\qquad
N(k)=\#\{(i,j)\in H\times(G\smallsetminus H)\mid i+j=k\}.
\]
\begin{itemize}[leftmargin=*]
\item $N(0)$: равенство $i+j=0$ означает $j=-i$, а $-i\notin H$ $\iff$ $-1\notin H$. Поэтому $N(0)=0$ при $p\equiv1\pmod4$ и
$N(0)=|H|=\frac{p-1}2$ при $p\equiv3\pmod4$.
\item \emph{$N(k)$ при $k\ne0$ не зависит от $k$.} Для $a\in G$ отображение $(i,j)\mapsto(ai,aj)$ при $a\in H$ и
$(i,j)\mapsto(aj,ai)$ при $a\notin H$ переводит $H\times(G\smallsetminus H)$ в себя (при $a\notin H$ имеем $aj\in H$,
$ai\notin H$) и решения уравнения $i+j=k$~--- в решения уравнения $i+j=ak$; обратное отображение строится так же по
$a^{-1}$. Значит, $N(ak)=N(k)$, а любое $k'\ne0$ имеет вид $ak$ с $a=k'k^{-1}$. Обозначим общее значение через $N$.
(Сдвиг $(i,j)\mapsto(i+c,j+c)$ здесь не годится: он не сохраняет $H$; например, при $p=7$ пара $(1,3)$ переходит в
$(2,4)$, а $4\in H$.)
\item Всего пар $|H|\cdot|G\smallsetminus H|=\bigl(\frac{p-1}2\bigr)^2$, т.\,е.\ $N(0)+(p-1)N=\frac{(p-1)^2}4$.
\end{itemize}
Так как $\sum_{k\ne0}\zeta^k=-1$, получаем $\alpha\beta=N(0)-N$, и
\begin{align*}
p\equiv1\pmod4:&\quad N=\tfrac{p-1}4,\quad \alpha\beta=-\tfrac{p-1}4,\quad X^2+X-\tfrac{p-1}4,\quad \alpha,\beta=\tfrac{-1\pm\sqrt p}2;\\
p\equiv3\pmod4:&\quad N=\tfrac{p-3}4,\quad \alpha\beta=\tfrac{p-1}2-\tfrac{p-3}4=\tfrac{p+1}4,\quad X^2+X+\tfrac{p+1}4,\quad
\alpha,\beta=\tfrac{-1\pm\sqrt{-p}}2 .
\end{align*}
В обоих случаях $(\alpha-\beta)^2=(\alpha+\beta)^2-4\alpha\beta=1-4\alpha\beta=p^*$, где $p^*=(-1)^{(p-1)/2}p$.

Пример: $p=7$, $H=\{1,2,4\}$:
\[
(\zeta+\zeta^2+\zeta^4)(\zeta^3+\zeta^5+\zeta^6)=\zeta^4+\zeta^6+\zeta^7+\zeta^5+\zeta^7+\zeta^8+\zeta^7+\zeta^9+\zeta^{10}
=3+(\zeta+\zeta^2+\dots+\zeta^6)=2=\tfrac{7+1}4 ;
\]
тройка~--- это $N(0)=3$ пары $(1,6)$, $(2,5)$, $(4,3)$ с суммой $0$, а каждая ненулевая степень $\zeta$ встречается $N=1$
раз. При $p=5$: $(\zeta+\zeta^4)(\zeta^2+\zeta^3)=\zeta^3+\zeta^4+\zeta+\zeta^2=-1=-\frac{5-1}4$; при $p=13$: $\alpha\beta=-3$.

\emph{Неподвижное поле.} По \rf{3.17}(b) $[E^H:\Q]=(G:H)=2$, и $\alpha\in E^H$ по (a). Число $\alpha=\frac{-1\pm\sqrt{p^*}}2$
иррационально, так как $p^*$ не является квадратом рационального числа. Значит, $\Q\subsetneq\Q(\alpha)\subset E^H$, откуда
$E^H=\Q(\alpha)=\Q(2\alpha+1)=\Q(\sqrt{p^*})$:
\[
E^H=\Q[\sqrt p]\ \ \text{при}\ p\equiv1\pmod4,\qquad E^H=\Q[\sqrt{-p}]\ \ \text{при}\ p\equiv3\pmod4 .
\]
Поскольку $H$~--- единственная подгруппа индекса $2$, это \emph{единственное квадратичное подполе} $E$ (\rf{3.17}). \qed

\begin{center}
\begin{keepfig}
\centering
\begin{tikzpicture}[every node/.style={font=\small}]
\node (Q) at (0,0) {$\Q$};
\node[cnew] (K) at (3.6,0) {$E^H=\Q\bigl(\sqrt{p^*}\bigr)=\Q(\alpha)$};
\node[cold] (E) at (7.6,0) {$E=\Q[\zeta]$};
\draw (Q)--(K)--(E);
\draw[decorate,decoration={brace,amplitude=4pt}] (0.1,0.3)--(3.0,0.3) node[midway,above=4pt] {$2$, \ $G/H\cong\Z/2\Z$};
\draw[decorate,decoration={brace,amplitude=4pt}] (4.9,0.3)--(7.5,0.3) node[midway,above=4pt] {$\frac{p-1}2$, \ $H$};
\draw[decorate,decoration={brace,amplitude=5pt,mirror}] (0.1,-0.3)--(7.5,-0.3) node[midway,below=5pt] {$p-1$, \ $G=\F_p^\times$};
\end{tikzpicture}
\figcap{Подгруппе квадратов $H$ отвечает квадратичное подполе; оно порождено гауссовым периодом $\alpha$.}
\end{keepfig}
\end{center}

\begin{remark}
(1) Разность $\alpha-\beta=\sum_{k=1}^{p-1}\bigl(\frac kp\bigr)\zeta^k$~--- это сумма Гаусса $g$ из \rf{5.11f}, где равенство
$g^2=p^*$ получено другим подсчётом; числа $\alpha$, $\beta$~--- гауссовы периоды подгруппы $H$ (\rf{5.11e}). (В \rf{5.11f}
и замечании после неё подгруппа $H$ обозначена $R$, а периоды $\alpha$, $\beta$~--- через $\eta_0$, $\eta_1$.)
Алгебраически знак в $\alpha=\frac{-1\pm\sqrt{p^*}}2$ не определён (автоморфизм из $G\smallsetminus H$ меняет $\alpha$ и
$\beta$ местами); при $\zeta=e^{2\pi i/p}$ всегда $\alpha=\frac{-1+\sqrt p}2$, соответственно $\alpha=\frac{-1+i\sqrt p}2$
(теорема Гаусса; см.\ текст после \rf{5.11f}).

(2) Вычислять $\alpha^2$ напрямую неудобно: в $\alpha^2=\sum_{h\in H}\zeta^{2h}+2\sum_{h<h'}\zeta^{h+h'}$ первая сумма равна
$\alpha$ или $\beta$ в зависимости от того, лежит ли $2$ в $H$, т.\,е.\ является ли $2$ квадратом по модулю $p$ (это так ровно
при $p\equiv\pm1\pmod8$~--- второе дополнение к квадратичному закону взаимности; например, $2\in H$ при $p=7,17$ и $2\notin H$
при $p=3,5,11,13$). Произведение $\alpha\beta$ этого знания не требует.

(3) Из упражнения следует, что каждое квадратичное поле $\Q(\sqrt d)$ содержится в некотором круговом поле (см.\ замечание
после \rf{5.11f}); общий результат такого рода~--- теорема Кронекера--Вебера: каждое конечное абелево расширение $\Q$
содержится в круговом поле.

(4) \emph{Орбиты смешанных пар.} Подсчёт $N(k)$~--- частный случай леммы \rf{3.22b}; наглядно он выглядит так. Группа
$G=\F_p^\times$ действует умножением на множестве $S$ неупорядоченных <<смешанных>> пар $\{i,j\}$, $i\in H$, $j\notin H$
(каждая такая пара однозначно записывается как $(i,j)\in H\times(G\smallsetminus H)$, так что
$\alpha\beta=\sum_{\{i,j\}\in S}\zeta^{i+j}$): при $a\in H$ вычет и невычет остаются вычетом и невычетом, при $a\notin H$
меняются ролями. Сумма пары умножается на $a$. Если $a\{i,j\}=\{i,j\}$, то либо $a=1$, либо $ai=j$, $aj=i$, т.\,е.\
$a^2=1$, $a=-1$ и $i+j=0$. Поэтому пары с суммой $0$ (они есть лишь при $p\equiv3\pmod4$) образуют одну орбиту из
$\frac{p-1}2$ пар $\{k,-k\}$, $k\in H$, со стабилизатором $\{\pm1\}$, а каждая пара с суммой $l\ne0$ имеет тривиальный
стабилизатор: её орбита состоит из $p-1$ пар, суммы которых $al$ пробегают все $\F_p^\times$ ровно по разу, и вклад орбиты в
$\alpha\beta$ равен $\sum_{l\ne0}\zeta^l=-1$. Итого при $p\equiv3\pmod4$~--- одна <<нулевая>> орбита и $\frac{p-3}4$
свободных (всего $\frac{p-1}2+\frac{p-3}4(p-1)=\bigl(\frac{p-1}2\bigr)^2$ пар), при $p\equiv1\pmod4$~--- $\frac{p-1}4$
свободных орбит; отсюда снова $\alpha\beta=\frac{p+1}4$, соответственно $-\frac{p-1}4$. Например, при $p=7$: орбита
$\{1,6\},\{2,5\},\{4,3\}$ и одна свободная орбита из шести пар; при $p=13$~--- три свободные орбиты по $12$ пар.
\end{remark}

% ---------------------------------------------------------------- 3-3
\begin{exercise}[3-3]
Пусть $M=\Q[\sqrt2,\sqrt3]$ и $E=M\bigl[\sqrt{(2+\sqrt2)(3+\sqrt3)}\bigr]$ (подполя $\R$).
\begin{enumerate}[label=(\alph*)]
\item Покажите, что $M$~--- расширение Галуа поля $\Q$ с группой Галуа $C_2\times C_2$ (четверная группа).
\item Покажите, что $E$~--- расширение Галуа поля $\Q$ с группой Галуа, изоморфной группе кватернионов $Q_8$.
\end{enumerate}
\end{exercise}

Напомним: $Q_8=\{\pm1,\pm i,\pm j,\pm k\}$ с умножением $i^2=j^2=k^2=ijk=-1$, где $-1$ коммутирует со всеми элементами.
Обозначим
\[
\gamma=(2+\sqrt2)(3+\sqrt3)=6+3\sqrt2+2\sqrt3+\sqrt6,\qquad \delta=\sqrt\gamma>0,
\]
так что $E=M[\delta]$. Поле $M$ и числа $\gamma$, $\delta$ уже изучены в \rf{1.32a}--\rf{1.32c} (там $M$ обозначалось $K$):
$1,\sqrt2,\sqrt3,\sqrt6$~--- базис $M$ над $\Q$; минимальный многочлен $\gamma$ равен $P=X^4-24X^3+144X^2-288X+144$ и
$M=\Q(\gamma)$; $\gamma$ не является квадратом в $M$, $[E:\Q]=8$, $E=\Q(\delta)$, а минимальный многочлен $\delta$ равен
$Q(X)=P(X^2)=X^8-24X^6+144X^4-288X^2+144$. (Минимальность $P$ теперь видна и иначе: автоморфизмы $M$, описанные в
(a) ниже, переводят $\gamma$ в четыре числа $6\pm3\sqrt2\pm2\sqrt3\pm\sqrt6$ (знак при $\sqrt6$ равен произведению двух
других), попарно различные по \rf{1.32a}; все они~--- корни минимального многочлена $\gamma$, так что его степень не меньше
$4=\deg P$.)

\medskip\noindent\emph{Решение.} (a) Поле $M$~--- поле разложения сепарабельного многочлена $(X^2-2)(X^2-3)$, поэтому $M/\Q$~---
расширение Галуа (\rf{3.10}), и $|\Gal(M/\Q)|=[M:\Q]=4$ (\rf{3.2}, \rf{1.32a}). Автоморфизм переводит корень $X^2-2$ в корень
$X^2-2$, поэтому $\sqrt2\mapsto\pm\sqrt2$, и так же $\sqrt3\mapsto\pm\sqrt3$; автоморфизм определяется этими двумя знаками
(так как $M=\Q[\sqrt2,\sqrt3]$). Комбинаций знаков четыре, автоморфизмов тоже четыре, значит, реализуются все:
\[
1,\qquad \sigma\colon\begin{cases}\sqrt2\mapsto-\sqrt2\\ \sqrt3\mapsto\sqrt3\end{cases},\qquad
\rho\colon\begin{cases}\sqrt2\mapsto\sqrt2\\ \sqrt3\mapsto-\sqrt3\end{cases},\qquad
\pi=\sigma\rho\colon\begin{cases}\sqrt2\mapsto-\sqrt2\\ \sqrt3\mapsto-\sqrt3\end{cases}.
\]
Квадрат каждого из них равен $1$, поэтому $\Gal(M/\Q)\cong C_2\times C_2$. (Отображение $\sqrt2\mapsto\sqrt3$ автоморфизмом быть
не может: соотношение $\sqrt2\cdot\sqrt2-2=0$ перешло бы в $\sqrt3\cdot\sqrt3-2=1\ne0$.) Заметим, что $\sigma$ и $\rho$ меняют
знак у $\sqrt6$, а $\pi$ его сохраняет. Неподвижные поля: $M^{\langle\sigma\rangle}=\Q(\sqrt3)$, $M^{\langle\rho\rangle}=\Q(\sqrt2)$,
$M^{\langle\pi\rangle}=\Q(\sqrt6)$ (каждое содержит указанное квадратичное поле и по \rf{3.17}(b) имеет степень $2$). Можно
рассуждать и так: группа порядка $4$~--- это $C_4$ или $C_2\times C_2$, но у $C_4$ одна подгруппа индекса $2$, а у $M$ три
различных квадратичных подполя.

Прежде чем переходить к (b), выясним, когда башня расширений Галуа сама является расширением Галуа.

\begin{lemma}[3.30e]
Пусть $F\subset M\subset E$~--- конечные расширения, причём $M/F$ и $E/M$~--- расширения Галуа.
\begin{enumerate}[label=\textup{(\alph*)}]
\item $E/F$~--- расширение Галуа тогда и только тогда, когда каждый $s\in\Gal(M/F)$ продолжается до автоморфизма поля
$E$. В этом случае ограничение $g\mapsto g|_M$~--- сюръективный гомоморфизм $\Gal(E/F)\to\Gal(M/F)$ с ядром $\Gal(E/M)$.
\item Пусть $\operatorname{char}F\ne2$ и $E=M[\sqrt a]$, где $a\in M$ не является квадратом в $M$. Автоморфизм
$s\in\Gal(M/F)$ продолжается до автоморфизма $E$ тогда и только тогда, когда $s(a)/a$~--- квадрат в $M$:
$s(a)=c_s^2\,a$, $c_s\in M$; продолжений ровно два: $\sqrt a\mapsto\pm c_s\sqrt a$. В частности, $E/F$~--- расширение
Галуа $\iff$ $s(a)/a\in M^{\times2}$ для всех $s\in\Gal(M/F)$.
\item В ситуации (b), если $E/F$~--- расширение Галуа и $s\in\Gal(M/F)$, $s\ne1$, $s^2=1$, то $c_s\,s(c_s)=\pm1$; оба
продолжения $s$ имеют порядок $2$, если $c_s\,s(c_s)=1$, и порядок $4$ (их квадрат~--- $\sqrt a\mapsto-\sqrt a$), если
$c_s\,s(c_s)=-1$.
\end{enumerate}
\end{lemma}

\begin{proof}
(a) Если $E/F$~--- расширение Галуа, то каждый $s\in\Gal(M/F)$, рассматриваемый как $F$-гомоморфизм $M\to E$,
продолжается до элемента $\Gal(E/F)$ по \rf{3.17b}. Обратно, пусть каждый $s$ продолжается. Любой $g\in\Aut(E/F)$
переводит $M$ в $M$: для $m\in M$ элемент $g(m)$~--- корень минимального многочлена $m$ над $F$, а он раскладывается в $M$
(расширение $M/F$ нормально). Поэтому $g|_M\in\Aut(M/F)=\Gal(M/F)$ (инъективное $F$-линейное отображение $M\to M$
биективно), и ограничение $r\colon\Aut(E/F)\to\Gal(M/F)$~--- гомоморфизм с ядром $\Aut(E/M)=\Gal(E/M)$, сюръективный по
предположению. Значит, $|\Aut(E/F)|=[E:M]\,[M:F]=[E:F]$ (\rf{3.2}, \rf{1.20}). По теореме Артина (\rf{3.11}) для
$G=\Aut(E/F)$ имеем $[E:E^G]=|G|=[E:F]$, а $F\subset E^G$; значит, $E^G=F$, и $E/F$~--- расширение Галуа (\rf{3.10}).
Утверждение об ограничении доказано попутно.

(b) Многочлен $X^2-a$~--- минимальный многочлен $\sqrt a$ над $M$, его корни $\pm\sqrt a$ различны ($\operatorname{char}\ne2$),
так что $E/M$~--- расширение Галуа. По \rf{2.4}(b) продолжения $s\colon M\to E$ до гомоморфизмов $E\to E$ взаимно
однозначно соответствуют корням многочлена $s(X^2-a)=X^2-s(a)$ в $E$; каждое такое продолжение~--- автоморфизм (инъективное
$F$-линейное отображение конечномерного пространства в себя). Корень в $E$ есть тогда и только тогда, когда $s(a)$~---
квадрат в $E$. По лемме из решения упражнения 2-1 элементы $M^\times$, являющиеся квадратами в $E=M[\sqrt a]$,~--- это
$M^{\times2}\cup aM^{\times2}$. Случай $s(a)\in M^{\times2}$ невозможен: тогда и $a=s^{-1}(s(a))$ был бы квадратом в $M$.
Итак, $s$ продолжается $\iff$ $s(a)=c_s^2a$ с $c_s\in M$, и тогда корни $X^2-s(a)$ в $E$~--- это $\pm c_s\sqrt a$.
Последнее утверждение следует из (a).

(c) Пусть $g$~--- продолжение $s$, $g(\sqrt a)=c\sqrt a$, $c=\pm c_s$. Тогда
\[
g^2(\sqrt a)=g(c)\,g(\sqrt a)=s(c)\,c\,\sqrt a=c_s\,s(c_s)\,\sqrt a
\]
(знак сокращается). Так как $g^2|_M=s^2=1$, то $g^2\in\Gal(E/M)=\{1,\ \sqrt a\mapsto-\sqrt a\}$, откуда $c_s\,s(c_s)=\pm1$, и
$g^2=1$ или $g^2\ne1$ соответственно; $g\ne1$, так как $g|_M=s\ne1$.
\end{proof}

Условие (b) можно сформулировать и так: \emph{$E/F$~--- расширение Галуа тогда и только тогда, когда в $E$ лежат все
корни многочлена $\prod_{s\in\Gal(M/F)}\bigl(X^2-s(a)\bigr)$} (его коэффициенты неподвижны относительно $\Gal(M/F)$, т.\,е.\
лежат в $F$). Действительно, по доказательству (b) автоморфизм $s$ продолжается на $E$ ровно тогда, когда $s(a)$~--- квадрат
в $E$, т.\,е.\ когда корни $\pm\sqrt{s(a)}$ многочлена $X^2-s(a)$ лежат в $E$. Именно так нормальность проверяется в шаге~3
решения ниже: все восемь чисел $\pm\sqrt{(2\pm\sqrt2)(3\pm\sqrt3)}$ лежат в $\Q(\delta)$.

\begin{example}
\begin{enumerate}[label=(\arabic*),leftmargin=*]
\item \emph{Башня расширений Галуа не обязана быть расширением Галуа} (ср.\ \rf{3.13a}(3)). Для башни
$\Q\subset\Q(\sqrt2)\subset\Q(\sqrt[4]2)=\Q(\sqrt2)[\sqrt a]$, $a=\sqrt2$, критерий (b) показывает это сразу:
$\sigma(a)/a=-1$ не является квадратом в $\Q(\sqrt2)\subset\R$. То же для $\Q\bigl(\sqrt{1+\sqrt2}\bigr)$:
$\frac{1-\sqrt2}{1+\sqrt2}=-(\sqrt2-1)^2$, а $-1$ не квадрат в $\Q(\sqrt2)$.
\item $\Q\bigl(\sqrt{2+\sqrt2}\bigr)$: $\frac{2-\sqrt2}{2+\sqrt2}=\frac{(2-\sqrt2)^2}2=(\sqrt2-1)^2$, $c_\sigma=\sqrt2-1$ и
$c_\sigma\sigma(c_\sigma)=(\sqrt2-1)(-\sqrt2-1)=-1$. По \rf{3.30e} это расширение Галуа степени $4$, и продолжение $\sigma$
имеет порядок $4$: группа циклическая, $\sqrt{2+\sqrt2}\mapsto(\sqrt2-1)\sqrt{2+\sqrt2}=\sqrt{2-\sqrt2}\mapsto-\sqrt{2+\sqrt2}\mapsto\dots$
\item Поле разложения $\Q(\sqrt[4]2,i)$ многочлена $X^4-2$ равно $M'[\sqrt a]$ при $M'=\Q(\sqrt2,i)$, $a=\sqrt2$ (здесь $\sqrt a=\sqrt[4]2\notin M'$, так как
$[\Q(\sqrt[4]2,i):\Q]=8$, см.\ пример после \rf{3.12a}). Для трёх
нетривиальных автоморфизмов $M'$ (меняющих знак у $\sqrt2$, у $i$ или у обоих) отношения $s(a)/a$ равны $-1=i^2$, $1$, $-1$;
$c_s=i,1,i$, а $c_s\,s(c_s)=i\cdot i=-1$, $1$, $i\cdot(-i)=1$. По \rf{3.30e}(c) получаем $2$ элемента порядка $4$ и
(вместе с $\sqrt a\mapsto-\sqrt a$) $5$ элементов порядка $2$: это диэдральная группа $D_4$ (см.\ отступление ниже).
\end{enumerate}
\end{example}

\noindent(b) \emph{Шаг 1: ключевые тождества.} Произведения $\gamma$ с его сопряжёнными над $\Q$ оказываются квадратами в $M$
(в каждом произведении одна из скобок <<сокращается>>):
\begin{align*}
\gamma\,\sigma(\gamma)&=(2+\sqrt2)(2-\sqrt2)(3+\sqrt3)^2=2(3+\sqrt3)^2=a_\sigma^2, & a_\sigma&=\sqrt2\,(3+\sqrt3),\\
\gamma\,\rho(\gamma)&=(2+\sqrt2)^2(3+\sqrt3)(3-\sqrt3)=6(2+\sqrt2)^2=a_\rho^2, & a_\rho&=\sqrt6\,(2+\sqrt2),\\
\gamma\,\pi(\gamma)&=(2+\sqrt2)(2-\sqrt2)(3+\sqrt3)(3-\sqrt3)=12=a_\pi^2, & a_\pi&=2\sqrt3 .
\end{align*}
Поэтому $s(\gamma)/\gamma=c_s^2$ с $c_s=a_s/\gamma\in M$. Явно (это удобнее находить прямо из отношений):
\[
\frac{\sigma(\gamma)}\gamma=\frac{2-\sqrt2}{2+\sqrt2}=\frac{(2-\sqrt2)^2}2=(\sqrt2-1)^2,\qquad
\frac{\rho(\gamma)}\gamma=\frac{3-\sqrt3}{3+\sqrt3}=\frac{(3-\sqrt3)^2}6=\Bigl(\frac{\sqrt6-\sqrt2}2\Bigr)^2,
\]
и $\pi(\gamma)/\gamma=\sigma\bigl(\rho(\gamma)/\gamma\bigr)\cdot\sigma(\gamma)/\gamma$. Итак,
\[
c_\sigma=\sqrt2-1,\qquad c_\rho=\frac{\sqrt6-\sqrt2}2,\qquad c_\pi=c_\sigma c_\rho=-1+\sqrt3+\frac{\sqrt2-\sqrt6}2
\]
(проверка: $c_\sigma\gamma=a_\sigma$, $c_\rho\gamma=a_\rho$, $c_\pi\gamma=a_\pi$).

\emph{Шаг 2: $E/\Q$~--- расширение Галуа, $[E:\Q]=8$.} Элемент $\gamma$ не является квадратом в $M$ (\rf{1.32c}). Это
можно увидеть и с помощью шага 1: если бы $\delta\in M$, то $\sigma(\delta)^2=\sigma(\gamma)=(c_\sigma\delta)^2$, т.\,е.\
$\sigma(\delta)=\varepsilon c_\sigma\delta$, $\varepsilon=\pm1$, и
\[
\delta=\sigma^2(\delta)=\sigma(\varepsilon c_\sigma\delta)=\varepsilon\sigma(c_\sigma)\cdot\varepsilon c_\sigma\delta=
(-\sqrt2-1)(\sqrt2-1)\,\delta=-\delta,
\]
что невозможно. (Ещё одно доказательство~--- через основную теорему: если $\gamma=\eta^2$, $\eta\in M$, то
$\bigl(\eta\,\sigma(\eta)\bigr)^2=\gamma\,\sigma(\gamma)=a_\sigma^2$, и $\sigma$-инвариантный элемент
$\eta\,\sigma(\eta)=\pm\sqrt2\,(3+\sqrt3)$ лежит в $M^{\langle\sigma\rangle}=\Q(\sqrt3)$, откуда $\sqrt2\in\Q(\sqrt3)$~---
противоречие с \rf{1.32a}.) По шагу 1 и \rf{3.30e}(b) поле $E=M[\delta]$~--- расширение Галуа $\Q$ степени $8$, а
$E=\Q(\delta)$ (\rf{1.32c}).

\emph{Шаг 3: восемь сопряжённых $\delta$ лежат в $\Q(\delta)$.} По \rf{3.30e}(b) группа $G=\Gal(E/\Q)$ состоит из
восьми автоморфизмов
\[
g_{s,\pm}\colon\ \delta\mapsto\pm c_s\delta,\qquad g_{s,\pm}|_M=s,\qquad s\in\{1,\sigma,\rho,\pi\},\ c_1=1 .
\]
Числа $\pm c_s\delta$~--- корни $Q$: $(c_s\delta)^2=s(\gamma)$ и $Q(\pm c_s\delta)=P(s(\gamma))=s(P(\gamma))=0$; они
попарно различны: из $c_s\delta=\pm c_t\delta$ следует $s(\gamma)=t(\gamma)$, т.\,е.\ $t^{-1}s$ тождествен на
$\Q(\gamma)=M$ (\rf{1.32b}), и $s=t$. Так как $c_s>0$ (например, $c_\pi\approx0{,}214$),
\begin{gather*}
\sqrt{(2-\sqrt2)(3+\sqrt3)}=(\sqrt2-1)\,\delta,\qquad \sqrt{(2+\sqrt2)(3-\sqrt3)}=\frac{\sqrt6-\sqrt2}2\,\delta,\\
\sqrt{(2-\sqrt2)(3-\sqrt3)}=c_\pi\,\delta :
\end{gather*}
все восемь корней $\pm\sqrt{(2\pm\sqrt2)(3\pm\sqrt3)}$ многочлена $Q$ выражаются через $\delta$, и $E=\Q(\delta)$~--- поле
разложения $Q$. Каждый элемент $G$ однозначно задаётся образом $\delta$ (\rf{2.1}(b)), и для каждого из восьми корней такой
элемент есть.

\emph{Шаг 4: порядки элементов.} Если $g(\delta)=c\,\delta$ и $h(\delta)=d\,\delta$ ($c,d\in M$), то
\[
gh(\delta)=g(d)\,g(\delta)=g|_M(d)\cdot c\cdot\delta .
\]
Автоморфизм $\tau=g_{1,-}\colon\delta\mapsto-\delta$ порождает $\Gal(E/M)$ и коммутирует со всеми $g$ (обе композиции
переводят $\delta$ в $-c\delta$). Для $s\ne1$ по \rf{3.30e}(c) порядок продолжений $s$ определяется знаком $c_s\,s(c_s)$:
\begin{align*}
c_\sigma\,\sigma(c_\sigma)&=(\sqrt2-1)(-\sqrt2-1)=-1,\\
c_\rho\,\rho(c_\rho)&=\tfrac{\sqrt6-\sqrt2}2\cdot\tfrac{-\sqrt6-\sqrt2}2=-\tfrac{6-2}4=-1,\\
c_\pi\,\pi(c_\pi)&=\bigl(c_\sigma\,\pi(c_\sigma)\bigr)\bigl(c_\rho\,\pi(c_\rho)\bigr)
=(\sqrt2-1)(-\sqrt2-1)\cdot\tfrac{\sqrt6-\sqrt2}2\cdot\tfrac{\sqrt6+\sqrt2}2=-1
\end{align*}
(здесь $\pi(\sqrt6)=\sqrt6$, $\pi(\sqrt2)=-\sqrt2$). Итак, \emph{все шесть элементов вне $\{1,\tau\}$ имеют порядок $4$, их
квадрат равен $\tau$, и $\tau$~--- единственный элемент порядка $2$ в $G$.} (Проверка знака здесь необходима для каждого $s$
по отдельности: <<симметрии>> между $\sigma$, $\rho$, $\pi$ нет, и в примере (3) выше знаки разные.)
Знак можно увидеть и без вычислений: так как $c_s=a_s/\gamma$ и $\gamma\,s(\gamma)=a_s^2$ (шаг 1),
\[
c_s\,s(c_s)=\frac{a_s\,s(a_s)}{\gamma\,s(\gamma)}=\frac{s(a_s)}{a_s},
\]
а $a_\sigma=\sqrt2\,(3+\sqrt3)$, $a_\rho=\sqrt6\,(2+\sqrt2)$, $a_\pi=2\sqrt3$, и автоморфизмы $\sigma$, $\rho$, $\pi$ меняют знак
соответственно у $\sqrt2$, $\sqrt6$, $\sqrt3$, сохраняя второй множитель. Иначе говоря, $\gamma\,s(\gamma)$ равно $2$, $6$, $3$,
умноженному на квадрат из $M$, и $s$ меняет знак у квадратного корня из этого числа~--- пример подобран именно так.

\emph{Шаг 5: $G\cong Q_8$.} Положим
\[
i=g_{\sigma,+}\colon\delta\mapsto c_\sigma\delta,\qquad j=g_{\rho,+}\colon\delta\mapsto c_\rho\delta,\qquad -1:=\tau,\qquad k=ij .
\]
По шагу 4 $i^2=j^2=\tau$; $k$ продолжает $\pi$, поэтому $k^2=\tau$; $ijk=(ij)^2=\tau$. По формуле шага 4
\[
ij(\delta)=\sigma(c_\rho)\,c_\sigma\delta=-c_\sigma c_\rho\delta=-c_\pi\delta,\qquad
ji(\delta)=\rho(c_\sigma)\,c_\rho\delta=c_\pi\delta,
\]
т.\,е.\ $ji=\tau\,ij$, и $G$ неабелева. Элементы $\pm1,\pm i,\pm j,\pm k$ группы $G$ (здесь $-g:=\tau g$) попарно различны
и удовлетворяют соотношениям $i^2=j^2=k^2=ijk=-1$; это соотношения группы кватернионов, и
$\pm1,\pm i,\pm j,\pm k\mapsto$ соответствующие элементы $G$~--- изоморфизм $Q_8\cong G$. (Иначе: $G$~--- неабелева группа
порядка $8$ с $j^2=i^2$ и $jij^{-1}=\tau i=i^{-1}$; по отступлению ниже это $Q_8$.) Сводка:
\begin{center}
\begin{tabular}{c|ccccc}
элемент $G$ & $1$ & $-1$ & $\pm i$ & $\pm j$ & $\pm k$\\\hline
$\delta\mapsto$ & $\delta$ & $-\delta$ & $\pm c_\sigma\delta$ & $\pm c_\rho\delta$ & $\mp c_\pi\delta$\\
ограничение на $M$ & $1$ & $1$ & $\sigma$ & $\rho$ & $\pi$\\
порядок & $1$ & $2$ & $4$ & $4$ & $4$
\end{tabular}
\end{center}


\emph{Шаг 6: подгруппы и подполя.} У $Q_8$ ровно шесть подгрупп: $\{1\}$, центр $\{\pm1\}$, три циклические подгруппы
$\langle i\rangle$, $\langle j\rangle$, $\langle k\rangle$ порядка $4$ и вся группа. Действительно, нетривиальная подгруппа
содержит $-1$ (элемент порядка $2$ и есть $-1$, а квадрат элемента порядка $4$ равен $-1$); подгруппа порядка $4$ содержит
элемент порядка $4$ (иначе в ней было бы три элемента порядка $2$) и потому циклическая, а циклических подгрупп порядка $4$
три: $\{\pm1,\pm i\}$, $\{\pm1,\pm j\}$, $\{\pm1,\pm k\}$. Соответствующие подполя: $E^{\{\pm1\}}=M$; так как
$-1\in\langle i\rangle$, поле $E^{\langle i\rangle}$ лежит в $M$ и состоит из элементов $M$, неподвижных относительно
$i|_M=\sigma$, т.\,е.\ $E^{\langle i\rangle}=M^{\langle\sigma\rangle}=\Q(\sqrt3)$; так же $E^{\langle j\rangle}=\Q(\sqrt2)$,
$E^{\langle k\rangle}=\Q(\sqrt6)$.

\begin{center}
\begin{keepfig}
\centering
\begin{tikzpicture}[every node/.style={font=\small}]
\begin{scope}
\node (e) at (0,3) {$\{1\}$};
\node (z) at (0,2) {$\{\pm1\}$};
\node (i) at (-1.5,1) {$\langle i\rangle$};
\node (j) at (0,1) {$\langle j\rangle$};
\node (k) at (1.5,1) {$\langle k\rangle$};
\node (G) at (0,0) {$Q_8$};
\draw (e)--(z) (z)--(i) (z)--(j) (z)--(k) (i)--(G) (j)--(G) (k)--(G);
\end{scope}
\node at (3.4,1.5) {$\longleftrightarrow$};
\begin{scope}[xshift=7.4cm]
\node[cnew] (L) at (0,3) {$E=\Q(\delta)$};
\node[cold] (K) at (0,2) {$M=\Q(\sqrt2,\sqrt3)$};
\node[cold] (a) at (-1.8,1) {$\Q(\sqrt3)$};
\node[cold] (b) at (0,1) {$\Q(\sqrt2)$};
\node[cold] (c) at (1.8,1) {$\Q(\sqrt6)$};
\node (Q) at (0,0) {$\Q$};
\draw (L)--(K) (K)--(a) (K)--(b) (K)--(c) (a)--(Q) (b)--(Q) (c)--(Q);
\draw[cnew] (-1.25,3) .. controls (-3.4,2.6) and (-3.4,0.3) .. (-0.35,0) node[pos=0.5,left] {$Q_8$};
\draw[cold] (1.55,2) .. controls (3.0,1.6) and (3.0,0.3) .. (0.35,0) node[pos=0.5,right] {$C_2\times C_2$};
\end{scope}
\end{tikzpicture}
\figcap{Решётка подгрупп $Q_8$ и решётка подполей $E$ (включения обращены). Все рёбра имеют индекс $2$; над каждым
квадратичным подполем поле $E$~--- циклическое расширение степени $4$.}
\end{keepfig}
\end{center}

\emph{Что здесь необычного.}
\begin{enumerate}[label=(\arabic*),leftmargin=*]
\item Все собственные подполя $E$ лежат в $M$ (каждая нетривиальная подгруппа содержит $-1$). Поэтому каждый элемент
$E\smallsetminus M$ порождает $E$ (например, $\delta+\sqrt2$ имеет степень $8$), $E$ не является композитом своих собственных
подполей, а подполе степени $4$ в $E$ единственно.
\item Все подгруппы $Q_8$ нормальны, поэтому \emph{каждое} подполе $E$~--- расширение Галуа над $\Q$ (\rf{3.17}(d)), хотя
группа $G$ неабелева. Для $D_4$ это не так: в поле разложения $X^4-2$ есть ненормальное подполе $\Q(\sqrt[4]2)$ (\rf{3.13a}(3), пример (1)
выше).
\item $E/\Q(\sqrt3)$~--- циклическое расширение степени $4$ с образующей $i$:
$\delta\mapsto c_\sigma\delta\mapsto-\delta\mapsto-c_\sigma\delta\mapsto\delta$, т.\,е.\
$\sqrt{(2+\sqrt2)(3+\sqrt3)}\mapsto\sqrt{(2-\sqrt2)(3+\sqrt3)}\mapsto-\sqrt{(2+\sqrt2)(3+\sqrt3)}\mapsto\dots$; так же над
$\Q(\sqrt2)$ и $\Q(\sqrt6)$.
\item $G/\{\pm1\}\cong\Gal(M/\Q)\cong C_2\times C_2$ (\rf{3.17}(d)), в согласии с (a). Но одно лишь наличие нормальной
подгруппы порядка $2$ с фактором $C_2\times C_2$ группу $G$ не определяет: решают порядки продолжений, т.\,е.\ знаки
$c_s\,s(c_s)$ (см.\ отступление). Какие биквадратные поля вкладываются в $Q_8$-расширение, описывает теорема Витта
(приводим без доказательства): при $\operatorname{char}F\ne2$ поле $F(\sqrt a,\sqrt b)$ степени $4$ содержится в
расширении Галуа поля $F$ с группой $Q_8$ тогда и только тогда, когда квадратичные формы $ax^2+by^2+abz^2$ и
$x^2+y^2+z^2$ эквивалентны над $F$. Для $a=2$, $b=3$ это так: при $x=\frac{v-w}2$, $y=\frac{u-v-w}3$, $z=\frac{2u+v+w}6$
(обратимая замена) $2x^2+3y^2+6z^2=u^2+v^2+w^2$. Например, $\Q(i,\sqrt2)$ в $Q_8$-расширение $\Q$ не вкладывается: форма
$-x^2+2y^2-2z^2$ принимает отрицательные значения. Это видно и непосредственно: если $\Q(i,\sqrt2)\subset L\subset\CC$ и
$\Gal(L/\Q)\cong Q_8$, то комплексное сопряжение задаёт элемент порядка $\le2$ группы $\Gal(L/\Q)$, т.\,е.\ $\pm1$, а оба
они тождественны на поле $\Q(i,\sqrt2)=L^{\{\pm1\}}$ (единственном подполе степени $4$, см.\ (1)), тогда как сопряжение
переводит $i$ в $-i$.
\end{enumerate}

\begin{aside}
\emph{Группы порядка $8$.} С точностью до изоморфизма их пять: абелевы $C_8$, $C_4\times C_2$, $C_2\times C_2\times C_2$ и
неабелевы $D_4$ (группа симметрий квадрата) и $Q_8$. Число элементов данного порядка:
\begin{center}
\begin{tabular}{lccccc}
группа & $C_8$ & $C_4\times C_2$ & $(C_2)^3$ & $D_4$ & $Q_8$\\\hline
порядка $2$ & $1$ & $3$ & $7$ & $5$ & $1$\\
порядка $4$ & $2$ & $4$ & $0$ & $2$ & $6$\\
порядка $8$ & $4$ & $0$ & $0$ & $0$ & $0$
\end{tabular}
\end{center}
Для неабелевых групп это проверяется так. Если все элементы имеют порядок $\le2$, то группа абелева ($ab=(ab)^{-1}=b^{-1}a^{-1}=ba$);
элемента порядка $8$ нет (иначе группа циклическая); значит, есть $a$ порядка $4$, и $A=\langle a\rangle$ имеет индекс $2$, т.\,е.\
нормальна. Для $b\notin A$ элемент $bab^{-1}\in A$ имеет порядок $4$, т.\,е.\ равен $a$ или $a^{-1}$, и в неабелевом случае
$bab^{-1}=a^{-1}$. Далее $b^2\in A$ (так как $G/A$ имеет порядок $2$) и $b^2$ коммутирует с $b$, а сопряжение с $b$ обращает
элементы $A$; поэтому $b^2=b^{-2}$, $b^2\in\{1,a^2\}$. При $b^2=1$ получаем $D_4$ ($a$~--- поворот на $90^\circ$, $b$~---
отражение), при $b^2=a^2$~--- $Q_8$ ($a=i$, $b=j$); в обоих случаях элементы однозначно записываются как $a^mb^n$
($0\le m<4$, $0\le n<2$), и соотношения определяют таблицу умножения.

Нормальная подгруппа $N$ порядка $2$ с фактором $C_2\times C_2$ есть во всех этих группах, кроме $C_8$ (фактор циклической группы цикличен, а $C_2\times C_2$ не
цикличен): в $C_4\times C_2$
это $\langle(2,0)\rangle$, в $(C_2)^3$~--- любая подгруппа порядка $2$, в $D_4$ и $Q_8$~--- центр $\{1,a^2\}$. В $D_4$ центр
порождён поворотом на $180^\circ$; он нормален (поворот на $180^\circ$ коммутирует со всеми симметриями), а остальные
четыре подгруппы порядка $2$, порождённые отражениями, не нормальны (сопряжение поворотом на $90^\circ$ переводит
отражение в другое). Фактор $D_4/\{1,a^2\}$ не имеет элементов порядка $4$ (квадрат любого элемента $D_4$ лежит в
$\{1,a^2\}$), т.\,е.\ тоже изоморфен $C_2\times C_2$. Так что по фактору группы не различить, а по числу элементов порядка
$2$~--- можно; именно это и сделано в шагах 4--5.
\end{aside}

\begin{aside}
\emph{Минимальные многочлены суммы и произведения.} Пусть $\beta,\gamma$ алгебраичны над $F$, $\beta_1,\dots,\beta_m$ и
$\gamma_1,\dots,\gamma_n$~--- корни их минимальных многочленов (в некотором поле разложения). Многочлены
\[
R_+(X)=\prod_{i,j}(X-\beta_i-\gamma_j)=\prod_i\mu_\gamma(X-\beta_i),\qquad R_\times(X)=\prod_{i,j}(X-\beta_i\gamma_j)
\]
имеют коэффициенты, симметричные отдельно по $\beta_i$ и по $\gamma_j$; по основной теореме о симметрических многочленах
(\rf{5.36}) это многочлены с целыми коэффициентами от коэффициентов $\mu_\beta$ и $\mu_\gamma$~--- явные формулы. Например,
для $\mu_\beta=X^2-sX+p$, $\mu_\gamma=X^2-tX+q$:
\[
R_\times=X^4-st\,X^3+(s^2q+t^2p-2pq)X^2-stpq\,X+p^2q^2 ;
\]
при $\beta=2+\sqrt2$ ($s=4$, $p=2$) и $\gamma=3+\sqrt3$ ($t=6$, $q=6$) получается многочлен $P$ из \rf{1.32b}, а для суммы
$5+\sqrt2+\sqrt3$ тем же способом~--- $R_+=X^4-20X^3+140X^2-400X+376$. Но минимальный многочлен $\beta+\gamma$ (или
$\beta\gamma$)~--- лишь неприводимый множитель $R_+$ (или $R_\times$), и $R$ может быть приводим: так бывает, когда числа
$\beta_i+\gamma_j$ не все различны или не все сопряжены (\rf{1.32d}(b): $\sqrt2+\sqrt8$; или $\beta=\sqrt2$, $\gamma=-\sqrt2$,
$R_+=X^2(X^2-8)$). Минимальность приходится доказывать отдельно, как в \rf{1.32b}; на языке теории Галуа минимальный
многочлен~--- это произведение $X-\theta$ по \emph{различным} сопряжённым $\theta$ элемента (\rf{3.14}).
\end{aside}

% ---------------------------------------------------------------- 3-4
\begin{exercise}[3-4]
Пусть $E$~--- расширение Галуа поля $F$ с группой Галуа $G$, и $L$~--- неподвижное поле подгруппы $H\subset G$. Покажите,
что группа автоморфизмов $L/F$ изоморфна $N/H$, где $N$~--- нормализатор $H$ в $G$.
\end{exercise}

Напомним, что нормализатор $N=N_G(H)=\{g\in G\mid gHg^{-1}=H\}$~--- это стабилизатор $H$ при действии $G$ на своих
подгруппах сопряжениями. Подгруппа $H$ нормальна в $N$, и $N$~--- наибольшая подгруппа $G$, в которой $H$ нормальна: если
$H\triangleleft K$, то каждый $k\in K$ нормализует $H$, т.\,е.\ $K\subset N$.

\medskip\noindent\emph{Идея.} Расширение $L/F$ не обязано быть расширением Галуа, и тогда автоморфизмов у $L$ меньше, чем $[L:F]$. Все
они получаются из элементов $G$: автоморфизм $L$ поднимается до автоморфизма $E$, и нужно лишь понять, какие $g\in G$
переводят $L$ в себя.

\medskip\noindent\emph{Решение.} \emph{Шаг 1: какие $g\in G$ сохраняют $L$.} По \rf{3.17}(c) $g(L)=g(E^H)=E^{gHg^{-1}}$, а
отображение $H'\mapsto E^{H'}$ инъективно (\rf{3.17}). Поэтому
\[
g(L)=L\iff gHg^{-1}=H\iff g\in N :
\]
стабилизатор подгруппы $H$ (относительно сопряжений) совпадает со стабилизатором поля $L$ (ср.\ \rf{3.17a}).

\emph{Шаг 2: гомоморфизм ограничения.} Для $g\in N$ ограничение $g|_L$~--- $F$-гомоморфизм $L\to L$, инъективный и
$F$-линейный, а значит, автоморфизм ($[L:F]<\infty$). Получаем гомоморфизм
\[
r\colon N\to\Aut(L/F),\qquad g\mapsto g|_L .
\]

\emph{Шаг 3: ядро.} $\ker r=\{g\in N\mid g|_L=\id\}=N\cap\Gal(E/L)=N\cap H=H$, так как $\Gal(E/E^H)=H$ (\rf{3.17}).

\emph{Шаг 4: сюръективность (подъём автоморфизмов).} Пусть $\varphi\in\Aut(L/F)$; рассматривая $\varphi$ как
$F$-гомоморфизм $L\to E$, продолжим его до некоторого $g\in G$ (\rf{3.17b}). Тогда $g(L)=\varphi(L)=L$, т.\,е.\ $g\in N$ по шагу 1, и $r(g)=\varphi$.

По теореме о гомоморфизме $N/H\cong\Aut(L/F)$, $gH\mapsto g|_L$. Иначе говоря, два элемента $N$ ограничиваются до одного
автоморфизма $L$ тогда и только тогда, когда они отличаются на элемент $H$, и подъёмы данного $\varphi$ образуют ровно
один смежный класс $gH$ (их $|H|=[E:L]$ штук). \qed

\begin{center}
\begin{keepfig}
\centering
\begin{tikzpicture}[every node/.style={font=\small}]
\node (F) at (0,0) {$F$};
\node[cnew] (K) at (2.5,0) {$E^N=L^{\Aut(L/F)}$};
\node[cold] (L) at (6.6,0) {$L=E^H$};
\node (E) at (9.3,0) {$E$};
\draw (F)--(K)--(L)--(E);
\draw[decorate,decoration={brace,amplitude=4pt},cnew] (3.9,0.3)--(6.5,0.3) node[midway,above=4pt] {$N/H\cong\Aut(L/F)$};
\draw[decorate,decoration={brace,amplitude=4pt}] (6.75,0.3)--(9.2,0.3) node[midway,above=4pt] {$H$};
\draw[decorate,decoration={brace,amplitude=4pt,mirror}] (3.9,-0.3)--(9.2,-0.3) node[midway,below=4pt] {$N$};
\draw[decorate,decoration={brace,amplitude=5pt,mirror}] (0.1,-0.9)--(9.2,-0.9) node[midway,below=5pt] {$G$};
\end{tikzpicture}
\figcap{Расширение $L/E^N$~--- расширение Галуа с группой $N/H\cong\Aut(L/F)$; все $F$-автоморфизмы $L$ тождественны на $E^N$.}
\end{keepfig}
\end{center}

\begin{remark}
(1) \emph{Подъём автоморфизмов.} Ключевой шаг~--- шаг 4: так как $E/F$ нормально, каждый $F$-гомоморфизм $L\to E$
продолжается до автоморфизма $E$ (\rf{3.17b}; для простого расширения $L=F[\alpha]$ это \rf{2.1}, \rf{2.4}: продолжения
соответствуют корням минимального многочлена). Без нормальности $E/F$ автоморфизм подняться не может: автоморфизм
$\sqrt2\mapsto-\sqrt2$ поля $\Q(\sqrt2)$ не продолжается на $\Q(\sqrt[4]2)$ (\rf{3.13a}(3), пример (1) после \rf{3.30e}).
Механизм подъёма такой: если $E=L[\beta]$ и $g$~--- минимальный многочлен $\beta$ над $L$, то продолжения
$\varphi\in\Aut(L/F)$ до гомоморфизмов $E\to E$ соответствуют корням многочлена $\varphi(g)$, а не $g$ (\rf{2.4}(b)), и в общем
случае $\varphi(g)\ne g$. Корни $\varphi(g)$ лежат в $E$ благодаря нормальности: $g$ делит минимальный многочлен $h$
элемента $\beta$ над $F$~--- многочлен с коэффициентами в $F$, так что $\varphi(h)=h$; поэтому $\varphi(g)$ делит $h$, а $h$
раскладывается в $E$ на линейные множители.

(2) \emph{Неподвижное поле автоморфизмов $L$.} Так как $\Aut(L/F)=r(N)$ и $E^N\subset E^H=L$,
\[
L^{\Aut(L/F)}=\{x\in L\mid g(x)=x\ \text{для всех}\ g\in N\}=E^N\cap L=E^N .
\]
По теореме Артина (\rf{3.11}) $L/E^N$~--- расширение Галуа с группой $N/H$ (это и \rf{3.17}(d) для $E/E^N$ и нормальной
подгруппы $H\triangleleft N$). Итак, $\Aut(L/F)=\Gal(L/E^N)$: каждый $F$-автоморфизм $L$ автоматически
тождествен на $E^N$. Более того, \emph{$E^N$~--- наименьшее промежуточное поле $F\subset k\subset L$, над которым
$L$~--- расширение Галуа}: если $L/k$~--- расширение Галуа, то $k=E^{H'}$ с $H'\supset H$, $E/k$~--- расширение Галуа с группой
$H'$ (\rf{3.13}), и по \rf{3.17}(d) для $E/k$ подгруппа $H$ нормальна в $H'$; значит, $H'\subset N$ и $k\supset E^N$. В
частности, $L/F$~--- расширение Галуа $\iff$ $N=G$ $\iff$ $H\triangleleft G$, в согласии с \rf{3.17}(d).

(3) \emph{Подсчёт.} $[L:F]=(G:H)=(G:N)(N:H)$, т.\,е.
\[
[L:F]=|\Aut(L/F)|\cdot[E^N:F],\qquad [E^N:F]=(G:N)=\text{число полей }gL,\ g\in G,
\]
сопряжённых с $L$ (\rf{3.17a}). В частности, $|\Aut(L/F)|$ делит $[L:F]$~--- это упражнение 3-5 для сепарабельных
расширений.

(4) \emph{Примеры.} Если $H\triangleleft G$, то $N=G$ и $\Aut(L/F)=G/H$. Для поля разложения $X^3-2$ над $\Q$ имеем
$G\cong S_3$, $L=\Q(\sqrt[3]2)$, $H$~--- подгруппа порядка $2$, $N=H$, и $\Aut(L/\Q)=1$ (\rf{3.3}); сопряжённых с $L$ полей три:
$3=1\cdot3$. Для $X^4-2$: $E=\Q(\theta,i)$, $\theta=\sqrt[4]2>0$, $[E:\Q]=8$ (пример после \rf{3.12a}); элементы $G$~--- это восемь отображений $\theta\mapsto i^m\theta$, $i\mapsto\pm i$ (каждое
действительно встречается, так как $|G|=8$). Возьмём $H=\{1,\kappa\}$, где $\kappa$~--- комплексное сопряжение; тогда
$L=E^H=\Q(\theta)$. Для $g\colon\theta\mapsto i^m\theta$, $i\mapsto\pm i$ имеем $g\kappa(\theta)=i^m\theta$ и
$\kappa g(\theta)=i^{-m}\theta$, так что $g$ коммутирует с $\kappa$ (т.\,е.\ $g\in N$) $\iff$ $m\in\{0,2\}$. Значит,
$|N|=4$, $\Aut(\Q(\theta)/\Q)\cong N/H\cong C_2=\{\id,\ \theta\mapsto-\theta\}$ (ср.\ \rf{3.3a}), $E^N=\Q(\sqrt2)$ (элементы $N$ сохраняют
$\theta^2=\sqrt2$), и $\Q(\theta)/\Q(\sqrt2)$~--- расширение Галуа (ср.\ пример (1) после \rf{3.30e}). Сопряжённых полей два:
$\Q(\theta)$ и $\Q(i\theta)$, и $4=2\cdot2$. Наконец, для $E$ из упражнения 3-3 все подгруппы нормальны, и
$\Aut(E^H/\Q)=G/H$ для всех $H$.
\end{remark}

% ---------------------------------------------------------------- 3-5
\begin{exercise}[3-5]
Пусть $E$~--- конечное расширение поля $F$. Покажите, что порядок группы $\Aut(E/F)$ делит степень $[E:F]$.
\end{exercise}

\noindent\emph{Решение.} Группа $G=\Aut(E/F)$ конечна: $|G|\le[E:F]$ по \rf{2.14}. Пусть $F'=E^G$; так как элементы $G$
тождественны на $F$, то $F\subset F'\subset E$. По теореме Артина (\rf{3.11}) $E/F'$~--- расширение Галуа с группой $G$ и
$[E:F']=|G|$. По \rf{1.20}
\[
[E:F]=[E:F']\,[F':F]=|G|\cdot[F':F],
\]
т.\,е.\ $|G|$ делит $[E:F]$. \qed

\begin{center}
\begin{keepfig}
\centering
\begin{tikzpicture}[every node/.style={font=\small}]
\node (F) at (0,0) {$F$};
\node[cnew] (K) at (3.4,0) {$F'=E^{\Aut(E/F)}$};
\node (E) at (6.6,0) {$E$};
\draw (F)--(K)--(E);
\draw[decorate,decoration={brace,amplitude=4pt}] (4.6,0.3)--(6.5,0.3) node[midway,above=4pt] {$|\Aut(E/F)|$, Галуа};
\draw[decorate,decoration={brace,amplitude=4pt}] (0.1,0.3)--(2.2,0.3) node[midway,above=4pt] {$[F':F]$};
\end{tikzpicture}
\figcap{Расширение $E$ всегда является расширением Галуа над неподвижным полем своей группы автоморфизмов.}
\end{keepfig}
\end{center}

Итак, \emph{любое} конечное расширение $E/F$ является расширением Галуа над $E^{\Aut(E/F)}$, и $\Aut(E/F)$~--- его группа
Галуа (\rf{3.5}); само $E/F$~--- расширение Галуа ровно тогда, когда $|\Aut(E/F)|=[E:F]$, т.\,е.\ $F'=F$ (\rf{3.10}).
Например, $|\Aut(\Q(\sqrt[4]2)/\Q)|=2$ (\rf{3.3a}), $F'=\Q(\sqrt2)$ (упражнение 3-4).

\begin{remark}
(1) \emph{Простые расширения.} Пусть $E=F[\alpha]$, $f$~--- минимальный многочлен $\alpha$. По \rf{3.3a}
$|\Aut(F[\alpha]/F)|$ равно числу $r$ различных корней $f$ в $F[\alpha]$, и по упражнению $r$ делит $\deg f$~--- в любой
характеристике, без предположений о сепарабельности или совершенности. Точнее, группа $\Aut(F[\alpha]/F)$ действует на этих
$r$ корнях свободно и транзитивно, и по доказательству \rf{3.10} минимальный многочлен $\alpha$ над $F'$ равен
$\prod(X-\alpha_i)$ по этим корням, так что $\deg f=r\cdot[F':F]$. Примеры: $\Q(\sqrt[3]2)$: $r=1$; $\Q(\sqrt[4]2)$: $r=2$
(корни $\pm\sqrt[4]2$); $\F_p(T)$ над $\F_p(T^p)$: $r=1$, $\deg f=p$ (\rf{3.3}); для расширения Галуа $r=\deg f$.

(2) \emph{Второе доказательство~--- через упражнение 3-4.} Если $E/F$ сепарабельно, то $E$ лежит в расширении Галуа
$K/F$ (\rf{3.12}); при $\Gamma=\Gal(K/F)$, $H=\Gal(K/E)$ имеем $E=K^H$ и по упражнению 3-4
\[
|\Aut(E/F)|=(N_\Gamma(H):H)\quad\text{делит}\quad(\Gamma:H)=[E:F].
\]
Без сепарабельности так рассуждать нельзя: минимальные многочлены образующих $E$ могут быть несепарабельны, и их поле
разложения не является расширением Галуа. Однако общий случай сводится к сепарабельному. Пусть $\operatorname{char}F=p>0$
и $E_s=\{x\in E\mid x\ \text{сепарабелен над}\ F\}$~--- подполе по \rf{2.22g}. Каждый $g\in\Aut(E/F)$ переводит $E_s$ на
себя ($g(x)$ имеет тот же минимальный многочлен, что и $x$). Если $g|_{E_s}=\id$, то $g=\id$: для $x\in E$ запишем его
минимальный многочлен как $h(X^{p^e})$ с наибольшим возможным $e$; тогда $h$ неприводим и не является многочленом от $X^p$,
т.\,е.\ сепарабелен (\rf{2.20}), так что $x^{p^e}\in E_s$; поэтому $g(x)^{p^e}=x^{p^e}$, $(g(x)-x)^{p^e}=0$ и $g(x)=x$. Значит,
ограничение $\Aut(E/F)\to\Aut(E_s/F)$ инъективно, и $|\Aut(E/F)|$ делит $|\Aut(E_s/F)|$ (теорема Лагранжа), которое делит
$[E_s:F]$ (сепарабельный случай), а $[E_s:F]$ делит $[E:F]$ (\rf{1.20}).
\end{remark}
